% Tracce uscite agli Esami di Stato (Sezione A, Settore Informazione) sugli argomenti del capitolo.
% Incluso da main.tex subito dopo il riassunto.

\section*{Tracce Uscite} % Usiamo * per non numerare la sezione nell'indice
\addcontentsline{toc}{section}{Tracce Uscite: Intelligenza Artificiale e ML}

\subsection*{Prima sessione 2024 -- Informatica \#3}
\textbf{Traccia}: tecniche di ricerca informata e non informata; vantaggi e svantaggi in termini di efficienza computazionale, memoria, ottimalità; esempi di problemi pratici risolvibili con una classe di tecniche rispetto all'altra.

\textbf{Dove trovare la risposta}: sezione \emph{Ricerca nello Spazio degli Stati}, con tabelle e scaletta di risposta.

\subsection*{Seconda sessione 2024 -- Informatica \#1}
\textbf{Traccia}: apprendimento supervisionato, non supervisionato e per rinforzo; principio e funzionamento; confronto su necessità di dati etichettati, uso dei dati, ambiti di applicazione; esempio concreto per ciascuno; limiti e casi in cui un approccio è più vantaggioso.

\textbf{Dove trovare la risposta}: sezione \emph{Paradigmi di Apprendimento}, con confronto in tabella e paragrafo \emph{Scegliere il paradigma}.

\subsection*{Prima sessione 2024 -- Automazione \#1}
\textbf{Traccia}: modellazione e analisi di sistemi dinamici; regressione lineare per identificazione e apprendimento automatico; underfitting e overfitting; scelta di training e validation set; metodo dei minimi quadrati con interpretazione geometrica, identificabilità e persistenza dell'eccitazione.

\textbf{Dove trovare la risposta}: sezioni \emph{Valutazione, Overfitting e Underfitting} e \emph{Regressione Lineare e Minimi Quadrati}.

\subsection*{Seconda sessione 2024 -- Automazione \#2}
\textbf{Traccia}: identificazione di sistemi dinamici con modelli parametrici lineari; predizione e simulazione; test di accuratezza e scelta dell'ordine; minimi quadrati e ricorsivi; identificazione non lineare con machine learning; generalizzazione, overfitting, underfitting, training e validation set.

\textbf{Dove trovare la risposta}: sottosezione \emph{Identificazione di sistemi dinamici} e sezioni su overfitting e minimi quadrati.

\subsection*{Prima sessione 2026 -- A3-INF-4 (tema di progetto)}
\textbf{Traccia}: nel sistema del cinema multisala è previsto un modulo di Machine Learning a supporto della pianificazione delle proiezioni.

\textbf{Dove trovare la risposta}: sezione \emph{Esempio: ML in un Sistema Completo}.
